%% ============================================================
%%  Lecture 7 -- Thursday, 18 June 2026
%%  Subfile of main.tex: compiles standalone OR as part of the book.
%%  Converted from math4061-current_notes/maths4061_lecture_07-2026_06_18/.
%% ============================================================
\documentclass[main.tex]{subfiles}

\begin{document}

\beginlecture{7}{Thursday, 18 June 2026}%
  {Function Limits, Continuity, and the Topology of Continuous Maps}

\epigraph{Natura non facit saltus.}{Gottfried Wilhelm Leibniz, New Essays IV.16}

\begin{lecturecontents}{Contents of this lecture}
  \item \hyperlink{7.1}{7.1\quad Function limits}
  \item \hyperlink{7.2}{7.2\quad Continuity and composition}
  \item \hyperlink{7.3}{7.3\quad Examples and discontinuities}
  \item \hyperlink{7.4}{7.4\quad Open-set characterizations}
  \item \hyperlink{7.5}{7.5\quad Continuous images of compact and connected sets}
  \item \hyperlink{7.6}{7.6\quad Homeomorphisms}
  \item \hyperlink{7.7}{7.7\quad Uniform continuity}
\end{lecturecontents}

\bigskip
\noindent\textbf{Overview.} This lecture begins Chapter~4 by replacing the sequential limit
$p_n\to p$ with the function limit $\lim_{x\to p}f(x)=q$ in metric spaces. The
$\eps$--$\delta$ definition is tied back to sequences, then specialized to continuity at a
point and continuity on a set. With the open-set characterization in hand, continuity becomes a
topological operation: preimages of open sets are open, compositions are continuous, compact
sets have compact images, connected sets have connected images, and the classical extreme-value
and intermediate-value theorems become consequences. The lecture closes with homeomorphisms and
the compactness theorem that upgrades pointwise continuity to uniform continuity.

\clearpage
% ============================================================
\sub{7.1}{Function limits}

\begin{Obj}{Theorem}{7.1.1}{Extreme values on a closed interval}
If $f:[a,b]\to\mathbb R$ is continuous, then $f$ attains a maximum and a minimum: there exist
$x,y\in[a,b]$ such that
\[
  f(x)\le f(z)\le f(y)\qquad\text{for every }z\in[a,b].
\]
\end{Obj}

\begin{Fig}{7.1.2}{The calculus picture of extrema}{A continuous real-valued function on a
closed interval cannot miss its highest and lowest values. The rigorous proof will come from
compactness of the image (\xref{7.5.2}).}
\begin{tikzpicture}[x=1cm,y=1cm]
  \draw[->] (-0.3,0) -- (7.2,0);
  \draw[->] (0,-0.3) -- (0,4.2);
  \draw[thick,figTerm]
    plot[smooth,tension=0.8] coordinates
    {(1,1.2) (2,3.4) (3,2.1) (4,2.7) (5,0.7) (6,2.0)};
  \draw[dashed] (1,0) -- (1,1.2);
  \draw[dashed] (6,0) -- (6,2.0);
  \node[below] at (1,-0.05) {$a$};
  \node[below] at (6,-0.05) {$b$};
  \filldraw[figLim] (2,3.4) circle (1.6pt);
  \draw[dashed,figLim] (2,3.4) -- (0,3.4);
  \draw[dashed,figLim] (2,3.4) -- (2,0);
  \node[below] at (2,-0.05) {$y$};
  \node[left,figLim] at (0,3.4) {$f(y)=\max$};
  \filldraw[figLim] (5,0.7) circle (1.6pt);
  \draw[dashed,figLim] (5,0.7) -- (0,0.7);
  \draw[dashed,figLim] (5,0.7) -- (5,0);
  \node[below] at (5,-0.05) {$x$};
  \node[left,figLim] at (0,0.7) {$f(x)=\min$};
\end{tikzpicture}
\end{Fig}

\begin{ObjL}{Remark}{7.1.3}{Sequences as functions}
For a sequence $\{p_n\}$ in a metric space $X$, writing $p_n\to p$ means
$\lim_{n\to\infty}p_n=p$ in the $\eps$--$N$ sense. A sequence is simply a function
\[
  f:\mathbb N\to X,\qquad n\mapsto p_n,
\]
and asks what happens as $n$ gets large.
\end{ObjL}

\begin{Obj}{Definition}{7.1.4}{Function limit in metric spaces}
Let $(X,d_X)$ and $(Y,d_Y)$ be metric spaces, let $E\subseteq X$, let $p$ be a cluster point of
$E$, and let $f:E\to Y$. We say that the \emph{limit of $f$ as $x\to p$ is $q$} if for every
$\eps>0$ there exists $\delta_{p,\eps}>0$ such that
\[
  0<d_X(x,p)<\delta_{p,\eps}
  \quad\Longrightarrow\quad
  d_Y\bigl(f(x),q\bigr)<\eps .
\]
We write
\[
  f(x)\to q\quad\text{as }x\to p,
  \qquad\text{or}\qquad
  \lim_{x\to p}f(x)=q .
\]
The condition $0<d_X(x,p)$ forces $x\ne p$; therefore $f(p)$ need not be defined, and even when
it is defined its value does not affect this limit.
\end{Obj}

\begin{Fig}{7.1.5}{The $\eps$--$\delta$ picture}{The punctured $\delta$-ball around $p$ in the
domain is carried inside the $\eps$-ball around $q$ in the target.}
\begin{tikzpicture}[x=1cm,y=1cm]
  \begin{scope}[shift={(0,0)}]
    \fill[figLim!12] (0,0) circle (1.3);
    \draw[figLim,thick] (0,0) circle (1.3);
    \draw[fill=white] (0,0) circle (0.06);
    \node[below right] at (0.02,0.02) {$p$};
    \filldraw[figTerm] (-0.5,0.5) circle (1.0pt);
    \filldraw[figTerm] (0.4,-0.4) circle (1.0pt);
    \filldraw[figTerm] (0.6,0.55) circle (1.0pt);
    \filldraw[figTerm] (-0.6,-0.45) circle (1.0pt);
    \draw[->,figLim] (0,0) -- (1.3,0);
    \node[above,figLim] at (0.65,0.0) {$\delta$};
    \node at (0,-1.75) {$(X,d_X)$};
  \end{scope}
  \draw[->,thick] (1.7,0.2) to[bend left=20] node[above]{$f$} (4.3,0.2);
  \begin{scope}[shift={(6,0)}]
    \draw[figLim,dashed,thick] (0,0) circle (1.3);
    \fill[figTerm!14]
      plot[smooth cycle,tension=0.9]
      coordinates {(-0.55,0.25) (-0.1,0.7) (0.45,0.45) (0.55,-0.1)
                   (0.2,-0.55) (-0.35,-0.4)};
    \filldraw[figLim] (0.15,-0.05) circle (1.2pt);
    \node[right,figLim] at (0.2,-0.05) {$q$};
    \draw[->,figLim] (0,0) -- (-0.92,0.92);
    \node[above left,figLim] at (-0.5,0.5) {$\eps$};
    \node at (0,-1.75) {$(Y,d_Y)$};
    \node[align=center] at (0,-2.35) {\small image of\\[-2pt]\small the $\delta$-ball};
  \end{scope}
\end{tikzpicture}
\end{Fig}

\begin{Obj}{Proposition}{7.1.6}{Sequential criterion for function limits}
Let $(X,d_X)$ and $(Y,d_Y)$ be metric spaces, let $E\subseteq X$, let $f:E\to Y$, and let $p$ be a
cluster point of $E$. Then
\[
  \lim_{x\to p}f(x)=q
  \quad\Longleftrightarrow\quad
  f(x_n)\to q
  \ \text{for every sequence } \{x_n\}\subseteq E\setminus\{p\} \text{ with } x_n\to p .
\]
\end{Obj}
\begin{Prf}{7.1.6}{Proof by the $\eps$--$\delta$ Definition and Its Contrapositive.}{[A2$+$A11]\quad use $\delta_n=1/n$ in the reverse direction.}
\item Suppose $\lim_{x\to p}f(x)=q$ and let $\{x_n\}\subseteq E\setminus\{p\}$ satisfy $x_n\to p$.
  Given $\eps>0$, choose $\delta>0$ such that
  \[
    0<d_X(x,p)<\delta\quad\Longrightarrow\quad d_Y(f(x),q)<\eps .
  \]
  Since $x_n\to p$, for all large $n$ one has $d_X(x_n,p)<\delta$, and the punctured condition is
  automatic because $x_n\ne p$. Hence $d_Y(f(x_n),q)<\eps$ eventually, so $f(x_n)\to q$.
\item Conversely, prove the contrapositive. If the function limit fails, then for some
  $\eps_0>0$ every $\delta>0$ fails. Taking $\delta=1/n$, choose $x_n\in E$ with
  \[
    0<d_X(x_n,p)<\frac1n
    \qquad\text{but}\qquad
    d_Y(f(x_n),q)\ge\eps_0 .
  \]
  Then $x_n\to p$, while $f(x_n)$ does not converge to $q$. This contradicts the sequential
  condition, so the function limit holds.
\end{Prf}

\noindent The logic is: if $A$ is the function-limit statement and $B$ is the sequential
statement, prove $A\Rightarrow B$ and $\neg A\Rightarrow\neg B$.

% ============================================================
\sub{7.2}{Continuity and composition}

\begin{Obj}{Definition}{7.2.1}{Continuity at a point}
Let $f:E\to Y$ and let $p\in E$. The function $f$ is \emph{continuous at $p$} if
\[
  \lim_{x\to p}f(x)=f(p).
\]
Equivalently, for every $\eps>0$ there exists $\delta_{\eps,p}>0$ such that
\[
  d_X(x,p)<\delta_{\eps,p}
  \quad\Longrightarrow\quad
  d_Y\bigl(f(x),f(p)\bigr)<\eps .
\]
For continuity the ball around $p$ is not punctured: the point $x=p$ is allowed, and the limit
value is required to be exactly $f(p)$.
\end{Obj}

\begin{Obj}{Definition}{7.2.2}{Continuity on a set}
A function $f:E\to Y$ is \emph{continuous on $E$} if it is continuous at every $p\in E$.
\end{Obj}

\begin{ObjL}{Remark}{7.2.3}{Isolated points}
If $p$ is an isolated point of $E$, then $f$ is continuous at $p$ vacuously: for a sufficiently
small ball, there are no points of $E$ near $p$ except $p$ itself.
\end{ObjL}

\begin{Obj}{Proposition}{7.2.4}{Composition of continuous functions}
If $f:X\to Y$ is continuous at $x$ and $g:Y\to Z$ is continuous at $f(x)$, then
$g\circ f:X\to Z$ is continuous at $x$. Consequently, the composition of continuous functions is
continuous.
\end{Obj}
\begin{Prf}{7.2.4}{Proof by an $\eps$--$\eta$--$\delta$ Chase.}{[A3]}
\item Let $\eps>0$. Since $g$ is continuous at $f(x)$, there exists $\eta>0$ such that
  \[
    d_Y(y,f(x))<\eta
    \quad\Longrightarrow\quad
    d_Z\bigl(g(y),g(f(x))\bigr)<\eps .
  \]
\item Since $f$ is continuous at $x$, there exists $\delta>0$ such that
  \[
    d_X(u,x)<\delta
    \quad\Longrightarrow\quad
    d_Y\bigl(f(u),f(x)\bigr)<\eta .
  \]
\item Combining the two implications, $d_X(u,x)<\delta$ implies
  \[
    d_Z\bigl((g\circ f)(u),(g\circ f)(x)\bigr)<\eps ,
  \]
  so $g\circ f$ is continuous at $x$.
\end{Prf}

\begin{Fig}{7.2.5}{The composition chase}{Continuity of $g$ supplies the $\eta$-ball in $Y$;
continuity of $f$ supplies the $\delta$-ball in $X$ whose image lands inside it.}
\begin{tikzpicture}[x=1cm,y=1cm]
  \begin{scope}[shift={(0,0)}]
    \draw[figLim,thick] (0,0) circle (0.95);
    \filldraw[figTerm] (0,0) circle (1.0pt) node[below right]{$x$};
    \draw[->,figLim] (0,0) -- (0.95,0);
    \node[above,figLim] at (0.5,0){$\delta$};
    \node at (0,-1.4){$X$};
  \end{scope}
  \draw[->,thick] (1.25,0.15) to[bend left=18] node[above]{$f$} (3.05,0.15);
  \begin{scope}[shift={(4.3,0)}]
    \draw[figLim,thick] (0,0) circle (0.8);
    \filldraw[figTerm] (0,0) circle (1.0pt) node[below right]{$f(x)$};
    \draw[->,figLim] (0,0) -- (0.8,0);
    \node[above,figLim] at (0.42,0){$\eta$};
    \node at (0,-1.4){$Y$};
  \end{scope}
  \draw[->,thick] (5.35,0.15) to[bend left=18] node[above]{$g$} (7.35,0.15);
  \begin{scope}[shift={(8.7,0)}]
    \draw[figLim,thick] (0,0) circle (1.05);
    \filldraw[figTerm] (0,0) circle (1.0pt) node[below right]{$g(f(x))$};
    \draw[->,figLim] (0,0) -- (-0.74,0.74);
    \node[above left,figLim] at (-0.4,0.4){$\eps$};
    \node at (0,-1.4){$Z$};
  \end{scope}
  \draw[->,thick,dashed] (0,-1.75) to[bend right=12]
        node[below]{$g\circ f$} (8.7,-1.75);
\end{tikzpicture}
\end{Fig}

\begin{ObjL}{Remark}{7.2.6}{Sequential proof of composition}
The sequential criterion gives the same result at a glance:
\[
  x_n\to x
  \quad\Longrightarrow\quad
  f(x_n)\to f(x)
  \quad\Longrightarrow\quad
  g(f(x_n))\to g(f(x)).
\]
\end{ObjL}

% ============================================================
\sub{7.3}{Examples and discontinuities}

\begin{ObjL}{Example}{7.3.1}{Basic continuous real functions}
The constant function $f(x)\equiv c$ is continuous: given $\eps>0$, every value of $f$ already
lies in $(c-\eps,c+\eps)$, so any $\delta>0$ works. The identity $f(x)=x$ is continuous by taking
$\delta=\eps$. Building from constants and the identity through the algebra of limits and
composition, every polynomial is continuous, and every rational function $p(x)/q(x)$ is
continuous where $q(x)\ne0$.
\end{ObjL}

\begin{Obj}{Proposition}{7.3.2}{The reciprocal is continuous on $(0,\infty)$}
The function $f(t)=1/t$ is continuous on $(0,\infty)$, but the required $\delta$ depends on both
the point $x$ and $\eps$.
\end{Obj}
\begin{Prf}{7.3.2}{Direct $\eps$--$\delta$ Proof for the Reciprocal.}{[A3]\quad the point-dependence is the point of the example.}
\item Fix $x>0$ and $\eps>0$. If $|t-x|<\delta$, then
  \[
    \left|\frac1t-\frac1x\right|=\frac{|t-x|}{tx}.
  \]
\item First keep $t$ away from $0$. If $\delta\le x/2$, then $t>x/2$, so
  $tx>x^2/2$ and
  \[
    \left|\frac1t-\frac1x\right|
    <\frac{\delta}{x^2/2}
    =\frac{2\delta}{x^2}.
  \]
\item Therefore it is enough to require $\delta<\eps x^2/2$. One may take
  \[
    \delta(x,\eps)=\min\!\left(\frac{x}{2},\,\frac{\eps x^2}{2}\right).
  \]
  This proves continuity at $x$. The factor $x^2$ shows why $\delta$ shrinks as $x\to0$.
\end{Prf}

\begin{Fig}{7.3.3}{Point-dependent $\delta$ for $1/x$}{The same $\eps$-band in the target pulls
back to a narrow input interval where the graph is steep and a wider interval where the graph is
flat.}
\begin{tikzpicture}[x=1cm,y=1cm]
  \draw[->] (-0.4,0) -- (5.8,0) node[right]{$x$};
  \draw[->] (0,-0.4) -- (0,3.45) node[above]{$y$};
  \draw[thick,figTerm]
        plot[smooth,variable=\x,domain=0.61:5.5,samples=160] ({\x},{2/\x});
  \node[figTerm,font=\small] at (4.55,0.85){$f(x)=\tfrac1x$};
  \filldraw[figTerm] (0.8,2.5) circle (1.3pt);
  \draw[dashed,figLim] (0,2.85) -- (0.70,2.85) -- (0.70,0);
  \draw[dashed,figLim] (0,2.15) -- (0.93,2.15) -- (0.93,0);
  \draw[<->,figLim,thick] (0.70,0.13) -- (0.93,0.13);
  \node[below,font=\scriptsize] at (0.8,-0.18){$b$};
  \node[figLim,font=\scriptsize] at (1.45,1.55){narrow $\delta$};
  \filldraw[figTerm] (2.5,0.8) circle (1.3pt);
  \draw[dashed,figLim] (0,1.15) -- (1.74,1.15) -- (1.74,0);
  \draw[dashed,figLim] (0,0.45) -- (4.44,0.45) -- (4.44,0);
  \draw[<->,figLim,thick] (1.74,0.13) -- (4.44,0.13);
  \node[below,font=\scriptsize] at (2.5,-0.18){$a$};
  \node[above,figLim,font=\scriptsize] at (3.05,0.18){wide $\delta$};
\end{tikzpicture}
\end{Fig}

\begin{ObjL}{Remark}{7.3.4}{Target and domain roles}
The $\eps$ always measures closeness in the target $Y$; the $\delta$ always measures closeness in
the domain $X$. In \xref{7.3.2} the roles do not swap. Only the size of $\delta$ changes with the
base point.
\end{ObjL}

\begin{ObjL}{Remark}{7.3.5}{Power series and familiar functions}
Power series define continuous functions inside their radius of convergence. In particular,
\[
  e^x=\sum_{n\ge0}\frac{x^n}{n!}
\]
is continuous, and similarly $\sin x$, $\cos x$, $\tan x$, and the other familiar functions are
continuous on their natural domains. Compare the radius-of-convergence discussion in
\xrefL{6.6.11}.
\end{ObjL}

\begin{ObjL}{Example}{7.3.6}{Standard discontinuities}
The floor function $\lfloor x\rfloor$ is discontinuous at each integer. The piecewise function
\[
  f(x)=
  \begin{cases}
    e^x, & x\ge0,\\
    0, & x<0,
  \end{cases}
\]
is discontinuous at $0$, since $\lim_{x\to0^+}f(x)=1$ while $\lim_{x\to0^-}f(x)=0$.
\end{ObjL}

\begin{Obj}{Definition}{7.3.7}{Type I and Type II discontinuities}
For a real function discontinuous at $x_0$, the one-sided limits classify the discontinuity.
\begin{itemize}[leftmargin=1.3em,topsep=2pt,itemsep=1pt,label=\textbullet]
  \item Type I: both one-sided limits exist, but either they differ (a jump) or they agree
        without matching $f(x_0)$ (a removable discontinuity).
  \item Type II: at least one one-sided limit fails to exist. This is often called an
        essential discontinuity.
\end{itemize}
\end{Obj}

\begin{Fig}{7.3.8}{Type I discontinuities}{A jump has different one-sided limits; a removable
discontinuity has a hole at the limiting value and the function value placed elsewhere.}
\begin{tikzpicture}[x=1cm,y=1cm]
  \begin{scope}
    \draw[->] (-0.2,0)--(3.2,0);
    \draw[->] (0,-0.3)--(0,2.6);
    \draw[thick,figTerm] (0.3,0.6) .. controls (1.0,0.9) .. (1.5,1.0);
    \draw[thick,figTerm] (1.5,1.9) .. controls (2.1,2.0) .. (2.8,2.1);
    \filldraw[figTerm] (1.5,1.0) circle (1.6pt);
    \draw[figTerm,fill=white] (1.5,1.9) circle (1.6pt);
    \draw[dashed] (1.5,0)--(1.5,1.0);
    \node[below,font=\scriptsize] at (1.5,-0.05){$x_0$};
    \node[font=\scriptsize] at (1.5,-0.7){jump};
  \end{scope}
  \begin{scope}[shift={(5,0)}]
    \draw[->] (-0.2,0)--(3.2,0);
    \draw[->] (0,-0.3)--(0,2.6);
    \draw[thick,figTerm] (0.3,1.1) .. controls (1.3,1.35) and (1.7,1.35) .. (2.8,1.1);
    \draw[figTerm,fill=white] (1.55,1.3) circle (1.6pt);
    \filldraw[figTerm] (1.55,2.0) circle (1.6pt);
    \draw[dashed] (1.55,0)--(1.55,1.3);
    \node[below,font=\scriptsize] at (1.55,-0.05){$x_0$};
    \node[font=\scriptsize] at (1.55,-0.7){removable};
  \end{scope}
\end{tikzpicture}
\end{Fig}

\begin{ObjL}{Example}{7.3.9}{The essential discontinuity of $\sin(1/x)$}
The function
\[
  f(x)=
  \begin{cases}
    \sin(1/x), & x\ne0,\\
    0, & x=0,
  \end{cases}
\]
has no limit at $0$. Near $0$ the graph oscillates infinitely often between $-1$ and $1$; for
different sequences $x_n\to0$, the values $f(x_n)$ can approach different limits.
\end{ObjL}

\begin{Fig}{7.3.10}{A Type II discontinuity}{The graph of $\sin(1/x)$ oscillates indefinitely
near $0$, so no single limiting value can govern all approaches to $0$.}
\begin{tikzpicture}[x=1cm,y=0.55cm]
  \draw[->] (-3.2,0)--(3.3,0) node[right]{$x$};
  \draw[->] (0,-1.5)--(0,1.6) node[above]{$y$};
  \draw[dashed,figLim] (-3.1,1)--(3.1,1) node[right,font=\tiny,black]{$1$};
  \draw[dashed,figLim] (-3.1,-1)--(3.1,-1) node[right,font=\tiny,black]{$-1$};
  \draw[thick,figTerm] plot[variable=\x,domain=0.07:3.1,samples=900]
        ({\x},{sin(deg(1/\x))});
  \draw[thick,figTerm] plot[variable=\x,domain=-3.1:-0.07,samples=900]
        ({\x},{sin(deg(1/\x))});
  \node[font=\scriptsize] at (0,-1.42){$f(x)=\sin(1/x)$ near $0$};
\end{tikzpicture}
\end{Fig}

% ============================================================
\sub{7.4}{Open-set characterizations}

\begin{Obj}{Proposition}{7.4.1}{Sequential characterization of continuity}
A function $f:E\to Y$ is continuous at $p\in E$ if and only if for every sequence
$\{x_n\}\subseteq E$ with $x_n\to p$, one has $f(x_n)\to f(p)$.
\end{Obj}
\begin{Prf}{7.4.1}{Proof by Specializing the Limit Criterion.}{[A2$+$A11]\quad the reverse direction again uses $\delta_n=1/n$.}
\item If $f$ is continuous at $p$, apply \xref{7.1.6} with $q=f(p)$; the non-punctured case causes
  no issue, since a sequence equal to $p$ along a subsequence already has image equal to $f(p)$
  along that subsequence.
\item Conversely, if continuity fails at $p$, then for some $\eps_0>0$ every $\delta>0$ fails.
  With $\delta=1/n$, choose $x_n\in E$ such that
  \[
    d_X(x_n,p)<\frac1n
    \qquad\text{but}\qquad
    d_Y(f(x_n),f(p))\ge\eps_0 .
  \]
  Then $x_n\to p$, but $f(x_n)$ does not converge to $f(p)$.
\end{Prf}

\begin{Obj}{Proposition}{7.4.2}{Continuity at a point via neighbourhoods}
The function $f:X\to Y$ is continuous at $p$ if and only if for every open set
$V\subseteq Y$ containing $f(p)$, there is an open set $U\subseteq X$ containing $p$ such that
\[
  f(U)\subseteq V
  \qquad\text{equivalently}\qquad
  U\subseteq f^{-1}(V).
\]
\end{Obj}
\begin{Prf}{7.4.2}{Proof by Translating Balls into Open Neighbourhoods.}{[A3]}
\item If $f$ is continuous at $p$ and $V\subseteq Y$ is open with $f(p)\in V$, choose
  $\eps>0$ such that $B_\eps(f(p))\subseteq V$. Continuity gives $\delta>0$ such that
  \[
    f(B_\delta(p))\subseteq B_\eps(f(p))\subseteq V.
  \]
  Taking $U=B_\delta(p)$ gives the required neighbourhood.
\item Conversely, take $V=B_\eps(f(p))$. The neighbourhood condition gives an open
  $U\ni p$ with $f(U)\subseteq V$. Since $U$ is open, some $B_\delta(p)\subseteq U$, and hence
  \[
    f(B_\delta(p))\subseteq B_\eps(f(p)).
  \]
  This is exactly the $\eps$--$\delta$ condition for continuity at $p$.
\end{Prf}

\begin{Obj}{Proposition}{7.4.3}{Global open-set characterization of continuity}
A function $f:X\to Y$ is continuous on $X$ if and only if $f^{-1}(V)$ is open in $X$ for every
open set $V\subseteq Y$.
\end{Obj}
\begin{Prf}{7.4.3}{Proof by Pulling Back Open Balls.}{[Direct]}
\item Suppose $f$ is continuous and let $V\subseteq Y$ be open. To show $f^{-1}(V)$ is open, take
  $x\in f^{-1}(V)$. Since $V$ is open, there is $\eps>0$ with
  $B_\eps(f(x))\subseteq V$. Continuity at $x$ gives $\delta>0$ with
  \[
    f(B_\delta(x))\subseteq B_\eps(f(x))\subseteq V.
  \]
  Hence $B_\delta(x)\subseteq f^{-1}(V)$, so every point of $f^{-1}(V)$ is interior.
\item Conversely, fix $x\in X$ and $\eps>0$. The ball $B_\eps(f(x))$ is open in $Y$, so
  $f^{-1}(B_\eps(f(x)))$ is open in $X$ by hypothesis. Since $x$ lies in this preimage, there is
  $\delta>0$ with
  \[
    B_\delta(x)\subseteq f^{-1}(B_\eps(f(x))).
  \]
  Thus $f(B_\delta(x))\subseteq B_\eps(f(x))$, so $f$ is continuous at $x$. Since $x$ was
  arbitrary, $f$ is continuous on $X$.
\end{Prf}

\begin{Obj}{Proposition}{7.4.4}{Composition via preimages}
The composition of continuous functions is continuous.
\end{Obj}
\begin{Prf}{7.4.4}{Proof by the Open-Set Characterization.}{[Direct]\quad a second proof of \xref{7.2.4}.}
\item Let $X\xrightarrow{\,f\,}Y\xrightarrow{\,g\,}Z$ with $f$ and $g$ continuous, and let
  $V\subseteq Z$ be open.
\item Since $g$ is continuous, $g^{-1}(V)$ is open in $Y$. Since $f$ is continuous,
  $f^{-1}(g^{-1}(V))$ is open in $X$.
\item But preimage commutes with composition:
  \[
    (g\circ f)^{-1}(V)=f^{-1}(g^{-1}(V)).
  \]
  Thus the preimage under $g\circ f$ of every open set is open, so $g\circ f$ is continuous.
\end{Prf}

\noindent This is the clean topological proof of composition: it works because preimages, not
images, commute exactly with composition.

\begin{ObjL}{Remark}{7.4.5}{Pugh's moral}
The sequential criterion is often the easiest way to check continuity at a glance: to test
continuity at $p$, push every sequence $x_n\to p$ through the function and see whether
$f(x_n)\to f(p)$.
\end{ObjL}

% ============================================================
\sub{7.5}{Continuous images of compact and connected sets}

\begin{Obj}{Theorem}{7.5.1}{Continuous image of a compact set}
If $f:X\to Y$ is continuous and $X$ is compact, then the image
\[
  f(X)=\{y\in Y:\exists\,x\in X\text{ with }f(x)=y\}
\]
is compact in $Y$.
\end{Obj}
\begin{Prf}{7.5.1}{Proof by Pulling Back an Open Cover.}{[A6]\quad uses \xref{7.4.3}.}
\item Let $\{V_\alpha\}$ be an open cover of $f(X)$, with each $V_\alpha$ open in $Y$.
  Since $f$ is continuous, each $f^{-1}(V_\alpha)$ is open in $X$.
\item These preimages cover $X$:
  \[
    X=\bigcup_\alpha f^{-1}(V_\alpha).
  \]
  Indeed, if $x\in X$, then $f(x)\in f(X)$, so $f(x)\in V_\alpha$ for some $\alpha$, and hence
  $x\in f^{-1}(V_\alpha)$.
\item By compactness of $X$, choose finitely many indices
  $\alpha_1,\dots,\alpha_n$ such that
  \[
    X=f^{-1}(V_{\alpha_1})\cup\cdots\cup f^{-1}(V_{\alpha_n}).
  \]
  Applying $f$, the sets $V_{\alpha_1},\dots,V_{\alpha_n}$ cover $f(X)$. Thus every open cover of
  $f(X)$ has a finite subcover, so $f(X)$ is compact.
\end{Prf}

\begin{Obj}{Corollary}{7.5.2}{Extreme value theorem}
If $f:[a,b]\to\mathbb R$ is continuous, then it attains a maximum and a minimum: there exist
$m,M\in[a,b]$ such that
\[
  f(m)\le f(x)\le f(M)\qquad\text{for all }x\in[a,b].
\]
\end{Obj}
\begin{Prf}{7.5.2}{Proof by Compactness of the Image.}{[A6$+$A5]\quad uses \xref{7.5.1} and Heine--Borel.}
\item The interval $[a,b]$ is compact, so $f([a,b])$ is compact by \xref{7.5.1}. By
  Heine--Borel in $\mathbb R$, it is closed and bounded.
\item Since $f([a,b])$ is bounded, the real numbers
  \[
    \alpha=\inf f([a,b]),
    \qquad
    \beta=\sup f([a,b])
  \]
  exist.
\item Since $f([a,b])$ is closed, it contains its infimum and supremum. Thus
  $\alpha,\beta\in f([a,b])$, so there are $m,M\in[a,b]$ with
  \[
    f(m)=\alpha,
    \qquad
    f(M)=\beta .
  \]
  Therefore $f(m)\le f(x)\le f(M)$ for every $x\in[a,b]$.
\end{Prf}

\begin{Obj}{Definition}{7.5.3}{Connected metric space}
A metric space $(X,d)$ is \emph{connected} if it cannot be written as
$X=A\cup B$ where
\[
  A\ne\emptyset,\qquad B\ne\emptyset,\qquad A\cap B=\emptyset,\qquad A,B\text{ are open}.
\]
Such a decomposition is called a \emph{separation}. Equivalently, the only clopen subsets of
$X$ are $\emptyset$ and $X$ itself. In $\mathbb R$, the connected subsets are exactly the
intervals.
\end{Obj}

\begin{Obj}{Theorem}{7.5.4}{Continuous image of a connected set}
If $X$ is connected and $f:X\to Y$ is continuous, then $f(X)$ is connected.
\end{Obj}
\begin{Prf}{7.5.4}{Proof by Pulling Back a Separation.}{[A9]\quad uses \xref{7.4.3}.}
\item Suppose for contradiction that $f(X)=A\cup B$ is a separation of $f(X)$: the sets
  $A$ and $B$ are nonempty, disjoint, and open in the subspace $f(X)$.
\item Then $f^{-1}(A)$ and $f^{-1}(B)$ cover $X$, are disjoint, and are open in $X$ by continuity.
  They are also nonempty: because $A,B\subseteq f(X)$ are nonempty, each contains some value
  $f(x)$.
\item Thus $X=f^{-1}(A)\cup f^{-1}(B)$ is a separation of $X$, contradicting connectedness.
  Therefore $f(X)$ is connected.
\end{Prf}

\begin{Obj}{Corollary}{7.5.5}{Intermediate value theorem}
Let $f:[a,b]\to\mathbb R$ be continuous and assume, without loss of generality, that
$f(a)\le f(b)$. Then for every $y\in[f(a),f(b)]$ there exists $x\in[a,b]$ with $f(x)=y$.
\end{Obj}
\begin{Prf}{7.5.5}{Proof by Connectedness of the Image.}{[Direct]\quad compactness gives closed and bounded; connectedness gives an interval.}
\item The interval $[a,b]$ is connected, so $f([a,b])$ is connected by \xref{7.5.4}. Since connected
  subsets of $\mathbb R$ are intervals, $f([a,b])$ is an interval.
\item This interval contains $f(a)$ and $f(b)$, and therefore contains every
  $y\in[f(a),f(b)]$. Hence for each such $y$ there is $x\in[a,b]$ with $f(x)=y$.
\end{Prf}

\noindent The structural mirror is useful: continuity preserves compactness, giving maxima and
minima; continuity preserves connectedness, giving intermediate values.

% ============================================================
\sub{7.6}{Homeomorphisms}

\noindent Two metric spaces are considered topologically the same when there is a bijection between
them that preserves the open-set structure in both directions.

\begin{Obj}{Definition}{7.6.1}{Homeomorphism}
A \emph{homeomorphism} is a continuous bijection $f:X\to Y$ whose inverse
$f^{-1}:Y\to X$ is also continuous. If such an $f$ exists, then $X$ and $Y$ are
\emph{homeomorphic}.
\end{Obj}

\begin{Obj}{Theorem}{7.6.2}{A continuous bijection from a compact domain is a homeomorphism}
Let $f:X\to Y$ be continuous and bijective. If $X$ is compact, then $f^{-1}:Y\to X$ is
continuous, so $f$ is a homeomorphism.
\end{Obj}
\begin{Prf}{7.6.2}{Proof by Showing the Inverse Pulls Back Opens to Opens.}{[A6]\quad compactness makes continuous images of closed sets closed.}
\item To prove $f^{-1}$ is continuous, it is enough to show that for every open
  $V\subseteq X$, the preimage of $V$ under $f^{-1}$ is open in $Y$. Since $f$ is a bijection,
  \[
    (f^{-1})^{-1}(V)=f(V).
  \]
\item Let $V\subseteq X$ be open. Then $V^c$ is closed in $X$, and since $X$ is compact,
  $V^c$ is compact.
\item By \xref{7.5.1}, $f(V^c)$ is compact in $Y$, hence closed. Because $f$ is a bijection,
  \[
    f(V)=\bigl(f(V^c)\bigr)^c .
  \]
  Therefore $f(V)$ is open in $Y$. Thus $f^{-1}$ is continuous.
\end{Prf}

\begin{Obj}{Counterexample}{7.6.3}{A continuous bijection need not be a homeomorphism}
Let
\[
  \varphi:[0,2\pi)\to S^1\subseteq\mathbb R^2,
  \qquad
  \varphi(t)=(\cos t,\sin t).
\]
This map is a continuous bijection onto the unit circle, but $[0,2\pi)$ is not compact and
$\varphi^{-1}$ is not continuous: points on $S^1$ approaching $(1,0)$ from the fourth quadrant
have arguments tending to $2\pi$, while $\varphi^{-1}(1,0)=0$.
\end{Obj}
\fails{the claim that every continuous bijection is automatically a homeomorphism.}

\begin{ObjL}{Remark}{7.6.4}{The compactness engine}
The proof of \xref{7.6.2} says that $f^{-1}$ is continuous exactly when $f$ carries open sets to
open sets; equivalently, it is enough here that $f$ carry closed sets to closed sets. On a compact
domain this is automatic: closed subsets of $X$ are compact, continuous images of compact sets are
compact, and compact subsets of a metric space are closed.
\end{ObjL}

\begin{ObjL}{Remark}{7.6.5}{Recognising homeomorphic spaces}
The basic homeomorphism questions are: given metric spaces $M$ and $N$, are they homeomorphic, and
what are the continuous maps $M\to N$? Typical comparison examples include the unit circle, a
trefoil knot, the perimeter of a square, the $2$-torus, $[0,1]$, and the unit disc.
\end{ObjL}

% ============================================================
\sub{7.7}{Uniform continuity}

\begin{Obj}{Definition}{7.7.1}{Uniform continuity}
A function $f:X\to Y$ is \emph{uniformly continuous} on $X$ if for every $\eps>0$ there exists
$\delta_\eps>0$ such that
\[
  d_X(p,q)<\delta_\eps
  \quad\Longrightarrow\quad
  d_Y\bigl(f(p),f(q)\bigr)<\eps .
\]
The difference from ordinary continuity is that $\delta$ depends only on $\eps$, not on the base
point.
\end{Obj}

\begin{Obj}{Theorem}{7.7.2}{Compactness gives uniform continuity}
If $X$ is compact and $f:X\to Y$ is continuous, then $f$ is uniformly continuous.
\end{Obj}
\begin{Prf}{7.7.2}{Proof by a Finite Half-Radius Cover.}{[A6$+$A4]\quad the half-radius leaves room for nearby points.}
\item Let $\eps>0$. Since $f$ is continuous at each $p\in X$, choose $\delta(p)>0$ such that
  \[
    d_X(p,q)<\delta(p)
    \quad\Longrightarrow\quad
    d_Y\bigl(f(p),f(q)\bigr)<\frac{\eps}{2}.
  \]
\item Define the half-radius ball
  \[
    J(p):=B^X_{\delta(p)/2}(p).
  \]
  The collection $\{J(p)\}_{p\in X}$ is an open cover of $X$, so compactness gives
  $p_1,\dots,p_N$ with
  \[
    X\subseteq J(p_1)\cup\cdots\cup J(p_N).
  \]
\item Set
  \[
    \delta=\frac12\min\bigl(\delta(p_1),\dots,\delta(p_N)\bigr)>0 .
  \]
  This $\delta$ depends only on $\eps$.
\item Take $x,y\in X$ with $d_X(x,y)<\delta$. Since the $J(p_i)$ cover $X$, choose $i$ with
  $x\in J(p_i)$. Then
  \[
    d_X(p_i,x)<\frac12\delta(p_i),
  \]
  so
  \[
    d_Y\bigl(f(p_i),f(x)\bigr)<\frac{\eps}{2}.
  \]
\item Also,
  \[
    d_X(y,p_i)\le d_X(y,x)+d_X(x,p_i)
    <\delta+\frac12\delta(p_i)
    \le\delta(p_i),
  \]
  so
  \[
    d_Y\bigl(f(y),f(p_i)\bigr)<\frac{\eps}{2}.
  \]
\item By the triangle inequality in $Y$,
  \[
    d_Y\bigl(f(y),f(x)\bigr)
    \le d_Y\bigl(f(y),f(p_i)\bigr)+d_Y\bigl(f(p_i),f(x)\bigr)
    <\frac{\eps}{2}+\frac{\eps}{2}
    =\eps .
  \]
  Thus $f$ is uniformly continuous.
\end{Prf}

\begin{Obj}{Counterexample}{7.7.3}{The reciprocal is not uniformly continuous on $(0,\infty)$}
The function $f(x)=1/x$ is continuous on $(0,\infty)$, but not uniformly continuous there. In the
explicit continuity estimate from \xref{7.3.2},
\[
  \delta(x,\eps)=\min\!\left(\frac{x}{2},\,\frac{\eps x^2}{2}\right)\to0
  \qquad\text{as }x\to0.
\]
Thus no single $\delta$ works for all $x>0$. The domain $(0,\infty)$ is not compact, so
\xref{7.7.2} does not apply. On a compact interval $[a,b]$ with $a>0$, the same function is
uniformly continuous.
\end{Obj}
\fails{the claim that ordinary continuity on a noncompact domain implies uniform continuity.}

\begin{ObjL}{Remark}{7.7.4}{Three compactness-for-free results}
Compactness upgrades continuity in three recurring ways: a continuous image of a compact set is
compact (\xref{7.5.1}); a continuous bijection from a compact domain is a homeomorphism
(\xref{7.6.2}); and a continuous map from a compact domain is uniformly continuous
(\xref{7.7.2}).
\end{ObjL}

% per-lecture Index of Objects -- standalone compile only; the book uses the master index.
\ifSubfilesClassLoaded{\clearpage\objindex}{}

\end{document}
